\section{Investigational Plan}

\subsection{Overall Study Design and Plan}

Study BREX-301 was a Phase III, randomized, double-blind, placebo-controlled, parallel-group, multicenter study. The overall study design is illustrated in Figure~\ref{fig:studydesign}.

\begin{figure}[H]
\centering
\begin{tikzpicture}[
  node distance=8mm and 2mm,
  phasebox/.style={rectangle, draw=accent, fill=accent!8, minimum width=28mm, minimum height=10mm, align=center, font=\footnotesize\bfseries, rounded corners=2pt},
  arrow/.style={-{Stealth[length=3mm, width=2mm]}, thick, accent},
]
\node[phasebox] (screen) {Screening\\$\le 4$ weeks};
\node[phasebox, right=of screen] (treat) {Treatment Period\\52 weeks};
\node[phasebox, right=of treat] (follow) {Follow-up\\4 weeks};
\node[above=2mm of screen, font=\footnotesize] {Visit 1};
\node[above=2mm of treat, font=\footnotesize] {Visit 2 (Baseline/Randomization)};
\node[above=2mm of follow, font=\footnotesize] {Week 52 (End of Treatment)};
\draw[arrow] (screen) -- (treat);
\draw[arrow] (treat) -- (follow);
\node[below=3mm of treat, font=\footnotesize, align=center] {\textbf{Randomization 2:1}\\[1pt]\studydrug\ \studydose\ QD \textit{vs.}\ Placebo QD};
\end{tikzpicture}
\caption{Overall Study Design}
\label{fig:studydesign}
\end{figure}

The study consisted of three periods: (1) a screening period of up to 4 weeks (Visits 1--2), during which patients were assessed for eligibility; (2) a 52-week double-blind treatment period (Visits 2--18), during which patients received either \studydrug\ \studydose\ or matching placebo once daily (QD); and (3) a 4-week post-treatment follow-up period (Visit 19). Study visits occurred at Weeks 0 (Baseline), 2, 4, 8, 12, 16, 24, 32, 40, and 52 during the treatment period.

\vspace{4pt}
Randomization was stratified by baseline weight category ($\le 90$ kg vs.\ $> 90$ kg) and prior biologic use for psoriasis (yes/no). An Interactive Response Technology (IRT) system was used to assign patients to treatment groups.

\subsection{Selection of Study Population}

\subsubsection{Inclusion Criteria}

Patients were eligible for enrollment if they met all of the following criteria:
\begin{enumerate}[label=\textbullet, leftmargin=*]
  \item Male or female, aged 18 to 75 years at the time of screening.
  \item Diagnosis of moderate-to-severe chronic plaque psoriasis for at least 6 months prior to screening.
  \item PASI score $\ge 12$ at both screening and baseline visits.
  \item sPGA score $\ge 3$ (moderate) at both screening and baseline visits.
  \item BSA involvement $\ge 10\%$ at both screening and baseline visits.
  \item Candidate for systemic therapy or phototherapy for psoriasis.
  \item Body weight $\ge 40$ kg and BMI $\le 40$ kg/m$^2$.
  \item Able to provide written informed consent.
\end{enumerate}

\subsubsection{Exclusion Criteria}

Key exclusion criteria included: non-plaque forms of psoriasis (guttate, erythrodermic, pustular); prior exposure to brexiprant; use of any biologic therapy within 5 half-lives or 12 weeks (whichever was longer) prior to baseline; use of systemic non-biologic psoriasis therapy within 4 weeks of baseline; active or latent tuberculosis; history of malignancy within 5 years (except treated non-melanoma skin cancer); clinically significant laboratory abnormalities at screening; pregnant or lactating women; and active infection requiring systemic antimicrobial therapy within 2 weeks of baseline.

\subsection{Treatments}

\subsubsection{Investigational Product}

Brexiprant was supplied as 300 mg film-coated tablets for oral administration. The dosage was one tablet taken orally once daily with water, with or without food, at approximately the same time each day. The 300 mg dose was selected based on the Phase II dose-ranging study (BREX-201), which demonstrated maximal efficacy at this dose with a favorable benefit--risk profile.

\subsubsection{Reference Therapy}

Matching placebo tablets, identical in appearance, size, and packaging to the active treatment, were administered orally once daily.

\subsubsection{Concomitant Therapy}

Low-to-moderate potency topical corticosteroids (Class III--VII) were permitted for treatment of the face, axillae, and groin throughout the study at the investigator's discretion. Emollients and moisturizers were allowed on all body areas except within 12 hours of a study visit. All other psoriasis therapies (topical, systemic, phototherapy) were prohibited during the study.

\subsection{Efficacy and Safety Variables}

The efficacy assessments included PASI scores, sPGA scores, DLQI scores, and BSA involvement at each study visit. Safety assessments included monitoring of treatment-emergent adverse events (TEAEs) using MedDRA version 27.0 coding, serious adverse events (SAEs), clinical laboratory evaluations (hematology, clinical chemistry, urinalysis), vital signs, 12-lead electrocardiograms, and physical examinations. TEAEs were defined as AEs with onset on or after the date of the first dose of study drug through 30 days after the last dose.

\subsection{Statistical Methods}

The Full Analysis Set (FAS) comprised all randomized patients who received at least one dose of study drug and had at least one post-baseline PASI assessment. The Per-Protocol (PP) analysis set included all FAS patients without major protocol deviations. The Safety Set included all randomized patients who received at least one dose of study drug.

\vspace{4pt}
The primary efficacy analysis compared the proportion of patients achieving PASI 75 at Week 16 between treatment groups using a Cochran--Mantel--Haenszel (CMH) test stratified by baseline weight category ($\le 90$ kg vs.\ $> 90$ kg) and prior biologic use (yes/no). Missing data were imputed using non-responder imputation (NRI) for the primary analysis. Sensitivity analyses using multiple imputation (MI) and tipping-point analyses were pre-specified.

\vspace{4pt}
A hierarchical testing procedure was applied to control the family-wise Type I error rate for the secondary endpoints, tested in the following pre-specified order: (1) sPGA 0/1 at Week 16, (2) PASI 90 at Week 16, (3) DLQI change from baseline at Week 16, and (4) PASI 75 at Week 52.

\begin{notebox}
The statistical analysis plan (SAP) was finalized on 15 July 2025, prior to database lock and unblinding. A cross-reference between the SAP and this CSR is provided in Section~\ref{sec:saper}.
\end{notebox}

\newpage

% ──────────────────────────────────────────────
% 6. STUDY PATIENTS
% ──────────────────────────────────────────────
